\documentclass[10pt,a4paper]{amsart}
\usepackage[margin=19mm]{geometry}
\usepackage{amsmath,amssymb,mathtools}
\usepackage[scr=rsfs,cal=euler]{mathalfa}
\usepackage{xfrac}
\usepackage[unicode,hidelinks,bookmarks=false]{hyperref}
\newcommand{\ie}{i.e.\ }
\newcommand{\cf}{cf.\ }
\newcommand{\resp}{respectively\ }
\newcommand{\Subsection}{\subsection}
\providecommand{\nobreakdash}{\nobreak\hbox{-}}
\setlength{\emergencystretch}{2em}
\multlinegap0pt
\usepackage{xcolor}
\usepackage{float}
\usepackage{mathtools}
\usepackage{amssymb}
\usepackage{tikz}
\usepackage[shortlabels]{enumitem}
\newtheorem{lemma}{Lemma}
\def\ZZ{{\mathbb Z}}
\newcommand{\ad}{\operatorname{ad}}
\def\gg{{\mathfrak{g}}}
\def\osp{{\mathfrak{osp}}}
\title{Perturbations of Dirac Operators}
\author{Steffen Schmidt}
\date{Source: arXiv:2603.23453v1 (CC BY 4.0)}
\begin{document}
\maketitle

\noindent\emph{Source and licence.} Perturbations of Dirac Operators. Version arXiv:2603.23453v1 (CC BY 4.0), distributed by arXiv under the Creative Commons Attribution 4.0 International licence, http://creativecommons.org/licenses/by/4.0/, as recorded in the arXiv licence metadata for this exact version and shown on the abstract page https://arxiv.org/abs/2603.23453v1. Full licence notice: \texttt{licenses/arxiv-2603.23453-CC-BY-4.0.txt}.\par\smallskip

\noindent\textit{Editorial scope.} Counted unit: the article's lemma lemm::properties\_phi\_g\_fd (source0.tex lines 1493--1504). The article's complete original proof of it is delivered with it (source0.tex lines 1506--1553, one \texttt{proof} environment of 48 lines). No mathematics is written, repaired or re-derived: the statement and the proof are the article's own lines, in the article's own notation, and no retained line is substituted or re-flowed.  Retained uncounted as the source-local context the proof needs: lines 1270--1272.  Nothing was omitted from any retained span, so the package carries no omission record.  No retained line contains a citation, so the wrapper carries no bibliography and no bibliography entry.  No retained line contains a figure environment, an image-inclusion command or drawing code, so no image is delivered and the package declares no asset dependency.  Editorial formatting only, in addition: an AMS \texttt{amsart} preamble that restates the article's own declarations and macros used by the retained lines, namely the environments \texttt{lemma} on separate counters, together with the standard packages those lines need.  The retained environments print in this document as Lemma 1 in order of appearance, whereas the article prints the same items with the numbering of its own full document.  The complete native source of the article as distributed in the arXiv e-print for this exact version is bundled unchanged as \texttt{sources/197088/source0.tex}.
\begin{lemma}\label{lemm::identification_lambda_omega}
Under $\gg\cong \gg^{\ast}$ one has $\lambda(T)=\omega_{T}$ for all $T\in\osp(\gg)$.
\end{lemma}

\begin{lemma} \label{lemm::properties_phi_g_fd} For all $x,y,z\in \gg$, the following hold:
 \begin{enumerate}
 \item[a)] $2\iota_{x}\phi=-\lambda(\ad_{x})$ for all $x \in \gg$.
 \item[b)] $\iota_{x}\iota_{y}\phi=B([x,y], \cdot)$. In particular, under the identification $\gg \cong \gg^{\ast}$, $\iota_{x}\iota_{y}\phi=[x,y]$.
 \item[c)] In the homogeneous basis, 
 \[
 \phi=-\frac{1}{12} \sum_{a,b,c} (-1)^{p(e_{a})p(e_{b})+p(e_{c})}f_{abc}e^{a}\wedge e^{b} \wedge e^{c}.
 \]
 \item[d)] $\phi$ is invariant under the adjoint action of $\gg$.
 
 \end{enumerate}
\end{lemma}

\begin{proof}
 a) 
For homogeneous $ x,y,z $, we obtain using Lemma~\ref{lemm::identification_lambda_omega}
\[
(\iota_x\phi)(y \wedge z)=\phi(x,y,z)
=-\tfrac{1}{2} B([x,y],z)
=-\tfrac{1}{2} \omega_{\ad_x}(y,z)=-\tfrac{1}{2}(\lambda(\ad_{x}))^{\flat}(y \wedge z).
\]
Hence $
2(\iota_x\phi)^\flat=-(\lambda(\ad_x))^\flat.
$
Using the $B$-identification $ \bigwedge^{2}(\gg) \cong \bigwedge^{2}(\gg^{\ast}) $, we finally obtain the statement.

 b) One has by definition:
 \[
 \begin{aligned}
 \iota_{x}\iota_{y}\phi(z) &=(\iota_{x}\iota_{y}\phi,z)_{\wedge}=(-1)^{p(x)p(y)}(\iota_{y}\phi, x \wedge z)_{\wedge}=(-1)^{p(x)p(y)}(\phi, y \wedge x \wedge z)_{\wedge} \\ &=-(-1)^{p(x)p(y)}B([y,x],z)=B([x,y],z).
 \end{aligned}
 \]

 c) is a straightforward computation and will be omitted. 

 d) One has $\phi(x,y,z)=-\tfrac12 B([x,y],z)$. For homogeneous $w,x,y,z$,
the Lie derivative on forms gives
\[
(L_w\phi)(x,y,z)
=\phi([w,x],y,z)+(-1)^{p(w)p(x)}\phi(x,[w,y],z)
+ (-1)^{p(w)(p(x)+p(y))}\phi(x,y,[w,z]).
\]
Hence, using invariance of $B$, one has for any $w \in \gg$
\[
\begin{aligned}
-2(L_w\phi)(x,y,z)
&=B([[w,x],y],z)
+ (-1)^{p(w)p(x)}B([x,[w,y]],z)
 +(-1)^{p(w)(p(x)+p(y))}B([x,y],[w,z])\\
&=B([w,x],[y,z])
 +(-1)^{p(w)p(x)}B(x,[[w,y],z])
 +(-1)^{p(w)(p(x)+p(y))}B(x,[y,[w,z]])\\
&=(-1)^{p(w)p(x)}\Big(
 -B(x,[w,[y,z]]) + B(x,[[w,y],z]) + (-1)^{p(w)p(y)}B(x,[y,[w,z]])
 \Big).
\end{aligned}
\]
By the $\ZZ_{2}$-graded Jacobi identity
$[w,[y,z]]=[[w,y],z] + (-1)^{p(w)p(y)}[y,[w,z]]$,
the parentheses vanish. Therefore $L_w\phi=0$ for all $w$, \emph{i.e.}, $\phi$ is invariant under the adjoint action of $\gg$.
\end{proof}

\end{document}
